% ============================================================
% Section 5: Red-Team Analysis
% ============================================================

\section{Red-Team Analysis}
\label{sec:redteam}

We adopt a deliberately adversarial posture toward our own claims.
This section reports the results of an external audit, a systematic seven-test red-team battery applied to the Phase~3d same-prompt data, the identification and characterization of the system prompt length confound, verification of the natural confabulation signal, and a transparent accounting of findings from earlier phases that proved to be artifacts.

% ============================================================
\subsection{External Adversarial Audit}
\label{sec:kavi_audit}

An external adversarial audit\footnote{Conducted by Kavi, an AI agent operating under an adversarial mandate. The audit independently executed 39 statistical tests across the codebase and experimental results.} evaluated 135 claims from the research program, identifying 14 discrepancies and conducting 39 independent statistical tests.
The audit was conducted with full access to all code, data, and result files.
We summarize the four most consequential findings:

\begin{enumerate}[leftmargin=2em, label=\textbf{D\arabic*.}]
  \item \textbf{Deception direction.}
  The original manuscript claimed that deception models showed a split between expansion and compression of \kvcache geometry.
  The audit verified that all seven deception models in Campaign~2 show expansion ($d < 0$ on spectral entropy, indicating richer geometry).
  The claimed compression split was an error in summarization; the underlying data were consistent throughout.

  \item \textbf{Bloom taxonomy inverted-U.}
  Campaign~2 reported an inverted-U relationship between Bloom cognitive complexity and geometric richness.
  The audit found this pattern is 90--98\% confounded with token count: higher-complexity prompts elicit longer responses, and the apparent inverted-U disappears after \fwl residualization.
  We retract this finding.

  \item \textbf{Deduplication.}
  The codebase included a \texttt{deduplicate\_runs()} function designed to remove duplicate trials, but this function was never called in any of the five experimental scripts that required it.
  Post-audit analysis confirmed that within-prompt identity similarity of 100\% (originally reported as a finding) was a data-leak artifact of non-deduplicated runs.
  After wiring deduplication into all affected scripts, cross-prompt identity similarity (92--97\%) remained robust, but within-prompt identity was no longer interpretable.

  \item \textbf{RDCT naming.}
  The ``Representational Distance under Cache Truncation'' (RDCT) paradigm was described as testing truncation of the \kvcache, but the actual experimental implementation applies \emph{prompt perturbation}---varying the prompt while measuring geometric change---not cache truncation.
  This invalidated a proposed connection to Watson's $1/e$ falsification test \cite{watson2024}, which requires genuine truncation.
  The experimental results themselves are valid; only the name and the claimed theoretical link were incorrect.
\end{enumerate}

All 14 discrepancies were resolved, with corrective commits merged via pull request to the main branch.
The deduplication fix was wired into all five affected experiment scripts; the deception narrative, Bloom finding, and RDCT description were corrected in all documentation.
We consider the audit process a model for how adversarial review should operate in interpretability research: full code access, independent reproduction, and public resolution.

% ============================================================
\subsection{Phase 3d Red-Team Battery}
\label{sec:redteam_battery}

Phase~3d applied seven systematic tests to the same-prompt data (150 trials per model: 50 prompts $\times$ 3 conditions $\times$ 3 models = 450 total trials).
Table~\ref{tab:redteam_summary} summarizes results; we detail each test below.

\begin{table}[t]
  \centering
  \caption{Phase~3d red-team battery summary.
  Tests T1--T3 and T5--T7 evaluate robustness of the detection signal.
  T4 (double residualization) tests for system prompt length confounds.
  T6 is informational only.
  Results reported as overall status across all three models (Qwen~7B, Llama~8B, Mistral~7B).}
  \label{tab:redteam_summary}
  \begin{tabular}{@{}clcl@{}}
    \toprule
    Test & Description & Status & Key Finding \\
    \midrule
    T1 & Permutation test (1{,}000 shuffles) & \textsc{pass} & $p = 0.000$ all comparisons \\
    T2 & Token-band analysis (terciles) & \textsc{pass} & Signal holds within all bands \\
    T3 & Encoding vs.\ delta decomposition & \textsc{pass} & Generation signal exceeds encoding \\
    T4 & Double residualization & \textsc{fail} & System prompt confound identified \\
    T5 & Behavioral compliance & \textsc{pass} & $\geq 94\%$ all conditions \\
    T6 & Demeaned analysis & \textsc{info} & Prompt-level signal strong \\
    T7 & Leave-one-prompt-out (LOPO) & \textsc{pass} & Generalizes to unseen prompts \\
    \bottomrule
  \end{tabular}
\end{table}

\paragraph{T1: Permutation test.}
The red-team battery uses all eight delta features ($\Delta$norm, $\Delta$norm\_per\_token, $\Delta$key\_rank, $\Delta$key\_entropy, $\Delta$top\_sv\_ratio, $\Delta$angular\_spread, $\Delta$norm\_variance, $\Delta$spectral\_entropy) rather than the four-feature classifier reported in Table~\ref{tab:sameprompt_auroc}; the resulting \auroc values are consequently higher (e.g., Qwen honest-vs-deceptive: 0.953 with eight features vs.\ 0.899 with four).
Condition labels were shuffled 1{,}000 times for each pairwise comparison.
None of the 1{,}000 permutations matched or exceeded the observed \auroc for any comparison on any model ($p = 0.000$ in all cases).
The observed \auroc values---Qwen 0.953, Llama 0.987, Mistral 0.982 for honest-vs-deceptive; Qwen 0.960, Llama 0.998, Mistral 1.000 for honest-vs-confabulation---exceed the permutation maxima by a wide margin.
On Qwen, the permutation maximum was 0.711 versus the observed 0.960 for honest-vs-confabulation; on Llama, the permutation maximum was 0.741 versus 0.987 for honest-vs-deceptive.
The signal is not attributable to label noise or chance partitioning.

\paragraph{T2: Token-band analysis.}
Response lengths were split into terciles (short, medium, long) and \auroc was recomputed within each band, controlling for gross length differences between conditions.
Table~\ref{tab:token_band} reports the mean within-band \auroc across all pairwise comparisons for each model.

\begin{table}[t]
  \centering
  \caption{Token-band analysis (T2): mean within-band \auroc across terciles (short, medium, long) for each model, averaged over all three pairwise comparisons.
  Signal persists within length-matched subsets, ruling out gross token-count confounds as the sole driver.}
  \label{tab:token_band}
  \begin{tabular}{@{}lcccc@{}}
    \toprule
    Model & Full \auroc & Short & Medium & Long \\
    \midrule
    Qwen 7B & 0.965 & 0.954 & 0.978 & 1.000 \\
    Llama 8B & 0.996 & 0.947 & 0.975 & 0.993 \\
    Mistral 7B & 0.999 & 0.996 & 0.995 & 0.986 \\
    \bottomrule
  \end{tabular}
\end{table}

\noindent Across all models, within-band \auroc exceeds 0.93 even in the shortest tercile, where token-count variation is most constrained.
The signal is robust to length stratification.

\paragraph{T3: Encoding vs.\ delta decomposition.}
This test separates the detection signal into encoding-phase and generation-phase components.
Table~\ref{tab:encoding_delta} reports the decomposition for the honest-vs-deceptive comparison, with \fwl residualization applied throughout.

\begin{table}[t]
  \centering
  \caption{Encoding vs.\ generation signal decomposition (T3, honest-vs-deceptive, eight-feature pipeline).
  Encoding \auroc reflects system prompt structure; generation \auroc reflects the model's response behavior; delta \auroc isolates the generation contribution.
  All values are \fwl-residualized.
  Note: the delta \fwl \auroc values here (e.g., Qwen 0.953) are from the eight-feature red-team pipeline and are higher than the four-feature values in Table~\ref{tab:sameprompt_auroc} (e.g., Qwen 0.899).}
  \label{tab:encoding_delta}
  \begin{tabular}{@{}lccc@{}}
    \toprule
    Model & Encoding \fwl \auroc & Generation \fwl \auroc & Delta \fwl \auroc \\
    \midrule
    Qwen 7B   & 0.819 & 0.974 & 0.953 \\
    Llama 8B  & 0.838 & 0.954 & 0.987 \\
    Mistral 7B & 0.720 & 0.966 & 0.982 \\
    \bottomrule
  \end{tabular}
\end{table}

\noindent Two findings emerge.
First, the generation-phase \fwl \auroc exceeds encoding-phase \fwl \auroc on every model, confirming that the detection signal originates from the model's own processing rather than from prompt structure alone.
Second, the encoding \fwl \auroc remains substantial (0.72--0.84), indicating that system prompt differences are partially visible in the key cache even after \fwl residualization.
This is expected: the different-length system prompts leave a structural fingerprint that standard token-count regression does not fully remove.
We address this directly in T4.

\paragraph{T4: Double residualization (\textsc{fail}).}
This is the test designed to expose the system prompt length confound.
When features are residualized against both system prompt length \emph{and} generated token count (Equation~\ref{eq:fwl-double}), the honest-vs-deceptive comparison weakens substantially:
\begin{equation*}
  \text{\auroc}_{\text{double-FWL}} =
  \begin{cases}
    0.726 & \text{(Qwen --- above chance)} \\
    0.750 & \text{(Llama --- above chance)} \\
    0.604 & \text{(Mistral --- near chance)}
  \end{cases}
\end{equation*}
Only Qwen (0.726) and Llama (0.750) retain honest-vs-deceptive signal above chance after double residualization; Mistral's 0.604 is near chance and is classified as a failure in the automated red-team pipeline.
All three models fail honest-vs-deceptive under the automated threshold; we interpret Qwen and Llama as retaining residual signal (above 0.70) while noting Mistral does not.
Comparisons involving forced confabulation collapse entirely.
Honest-vs-confabulation \auroc drops to 0.500 (Qwen), 0.411 (Llama), and 0.531 (Mistral)---at or below chance.
Confabulation-vs-deceptive \auroc drops to 0.532 (Qwen), 0.568 (Llama), and 0.710 (Mistral).

The diagnosis is clear: the Phase~3d forced confabulation system prompt differs in length from the honest and deceptive system prompts, creating a system prompt length confound that drives the confabulation signal.
When that confound is removed via double residualization, the confabulation comparisons collapse while the honest-vs-deceptive comparison---where the system prompt length difference is smaller---retains residual signal.

This failure is instructive rather than fatal.
The Phase~2b experiments (Section~\ref{sec:misalignment_detection}) used equal-length padded system prompts, and the honest-vs-deceptive detection in that design achieves \fwl \auroc of 0.935--0.995, which is not vulnerable to this confound.
The Phase~3d double-FWL honest-vs-deceptive values of 0.73--0.75 (on Qwen and Llama) represent a stringent lower bound: the genuine generation-phase signal that survives after removing both token count and system prompt length effects.

\paragraph{T5: Behavioral compliance.}
Compliance verification confirmed that models followed system prompt instructions in the overwhelming majority of trials: compliance rates ranged from 94\% (Llama forced confabulation) to 100\% (several conditions across all models).
Hedge rates were low (0--6\% across misaligned conditions), confirming that the models committed to the instructed behavior rather than equivocating.
These compliance rates validate that the geometric signals reflect genuine condition differences rather than models partially ignoring instructions.

\paragraph{T6: Demeaned analysis (\textsc{info}).}
Prompt-level demeaning---subtracting each prompt's mean across conditions before classification---isolates variation \emph{within} prompt-condition pairs from variation \emph{across} prompts.
For honest-vs-deceptive, demeaned \fwl \auroc matched or exceeded raw \fwl \auroc (e.g., Qwen: raw 0.953, demeaned 1.000; Llama: raw 0.987, demeaned 1.000), indicating that the discriminative signal is present at the within-prompt level, not driven by between-prompt variation.
\textbf{Selective reporting caveat:} We report demeaned results here only for honest-vs-deceptive.
Demeaned analysis was not systematically applied to confabulation comparisons (honest-vs-confabulation, confabulation-vs-deceptive).
Because confabulation comparisons collapse under double residualization (T4), demeaned results for those comparisons would be difficult to interpret and could create a misleading impression of robustness.
We flag this gap for transparency: future work should apply demeaned analysis uniformly across all comparisons and report results regardless of whether they are favorable.
This test is informational because demeaning can introduce overfitting when the number of conditions per prompt is small (here, 3).

\paragraph{T7: Leave-one-prompt-out (LOPO).}
A classifier trained on 49 prompts and tested on the held-out prompt was iterated across all 50 prompts.
Table~\ref{tab:lopo} reports LOPO performance.

\begin{table}[t]
  \centering
  \caption{Leave-one-prompt-out (T7) results.
  LOPO \auroc closely tracks the full-data \auroc, demonstrating generalization to unseen prompt content.}
  \label{tab:lopo}
  \begin{tabular}{@{}lcccccc@{}}
    \toprule
    & \multicolumn{2}{c}{Qwen 7B} & \multicolumn{2}{c}{Llama 8B} & \multicolumn{2}{c}{Mistral 7B} \\
    \cmidrule(lr){2-3} \cmidrule(lr){4-5} \cmidrule(lr){6-7}
    Comparison & Full & LOPO & Full & LOPO & Full & LOPO \\
    \midrule
    H vs.\ D    & 0.953 & 0.954 & 0.987 & 0.978 & 0.982 & 0.980 \\
    H vs.\ C    & 0.960 & 0.956 & 0.998 & 1.000 & 1.000 & 1.000 \\
    C vs.\ D    & 0.915 & 0.930 & 0.965 & 0.962 & 0.995 & 0.996 \\
    \midrule
    3-way       & 0.965 & 0.963 & 0.996 & 0.999 & 0.999 & 0.999 \\
    \bottomrule
  \end{tabular}
\end{table}

\noindent LOPO \auroc closely tracks full-data \auroc across all comparisons and models, with several LOPO values actually exceeding the full-data value (e.g., Mistral C-vs-D: full 0.995, LOPO 0.996).
This confirms that the classifier generalizes to unseen prompt content and is not memorizing prompt-specific patterns.
LOPO accuracy exceeds 87\% in all comparisons (ranging from 87\% for Qwen C-vs-D to 100\% for Mistral H-vs-C).

% ============================================================
\subsection{The System Prompt Length Confound}
\label{sec:prompt_length_confound}

The T4 failure motivates a detailed discussion of the system prompt length confound, which is the single most important methodological lesson of this research program.

\paragraph{The mechanism.}
When system prompts differ in token count, the encoding-phase key cache has different dimensionality across conditions before generation even begins.
A difference as small as 14--15 tokens produces encoding \auroc of 0.744 (Phase~2b, before padding was introduced).
This occurs because the \kvcache stores one key-value pair per token at each layer: a longer system prompt creates a physically larger cache, which affects every downstream feature (norm, rank, entropy) through the dependence on sequence length $T$ in Equations~\ref{eq:kflat}--\ref{eq:norm-per-token}.

Standard \fwl residualization against total token count $T$ mitigates but does not fully remove this confound, because the system prompt tokens contribute to $T$ in a fixed, condition-specific manner that is collinear with the condition label.
Only double residualization (Equation~\ref{eq:fwl-double}), which explicitly models the system prompt contribution, can isolate the generation signal from the prompt-length artifact.

\paragraph{The solution.}
The solution is simple: \textbf{pad all system prompts to equal token counts} before tokenization, then verify post-tokenization equality.
Phase~2b adopted this protocol, and all subsequently reported \auroc values from Phase~2b (0.935--0.995) are computed under equal-length conditions.
Under equal-length padding, encoding \auroc drops to near chance, confirming that the prompt structure confound has been eliminated.

\paragraph{Bounding the residual signal.}
For Phase~3d, where system prompts were not padded, double residualization provides a conservative lower bound on the genuine detection signal.
The honest-vs-deceptive \auroc of 0.73--0.75 (Qwen and Llama) under double residualization represents the signal that survives after controlling for \emph{both} token count and system prompt length.
We interpret this as the floor of the genuine cognitive geometry signal, with the Phase~2b equal-length \auroc of 0.935--0.995 representing the achievable performance when the confound is eliminated by design rather than by statistical correction.

% ============================================================
\subsection{Natural Confabulation}
\label{sec:natural_confab}

Phase~3c tested whether confabulation detection generalizes beyond forced conditions to \emph{natural} confabulation---instances where a model spontaneously generates false information when asked about fictitious entities.
We presented each model with 25 questions about nonexistent people, places, equations, and events (e.g., ``What was the Treaty of Kensington signed in 1743?'') alongside 25 factual baseline questions with verifiable answers (e.g., ``What is the chemical symbol for gold?'').

\paragraph{Confabulation rates.}
Confabulation rates varied substantially across models.
Qwen confabulated on 21/25 fictitious prompts, Mistral on 23/25, and Llama on only 15/25 (with 9 hedged responses of the form ``I couldn't find any information\ldots'').
This model-level variation in confabulation propensity is itself informative: Llama appears more calibrated to its knowledge boundaries than Qwen or Mistral, consistent with the hedging behavior expected from RLHF-aligned models with stronger refusal tuning.

\paragraph{Detection after FWL.}
Raw \auroc for factual-vs-confabulated was near-ceiling across all models (Qwen 0.997, Llama 0.950, Mistral 0.997).
However, after \fwl residualization against token count, detection drops substantially:
\begin{equation*}
  \text{\auroc}_{\text{FWL}} =
  \begin{cases}
    0.704 & \text{(Qwen)} \\
    0.624 & \text{(Llama)} \\
    0.725 & \text{(Mistral)}
  \end{cases}
\end{equation*}
Encoding \auroc remains high (0.957--0.970), indicating that the different prompt categories (fictitious questions vs.\ factual questions) are distinguishable from their content structure alone, before any generation occurs.

\paragraph{Verdict.}
The natural confabulation signal is largely a \textbf{content confound}: fictitious prompts and factual prompts differ systematically in structure, and the raw detection performance reflects this structural difference rather than (or in addition to) the model's cognitive state during confabulation.
After \fwl control, the residual signal (0.62--0.73) is modest and is not cleanly separable from prompt structure effects.
We classify Qwen as ``confound-controlled'' (weak but surviving signal) and Llama and Mistral as ``null result'' (residual signal does not clearly exceed what prompt structure alone can explain).

This negative result reinforces a key methodological point: \textbf{same-prompt designs are essential} for clean detection of cognitive mode.
The Phase~3d forced-confabulation results (where the same prompt is answered honestly or confabulatorily depending on the system prompt) achieve \auroc of 0.92--1.00, compared to the 0.62--0.73 of the cross-prompt Phase~3c design.
The difference is attributable almost entirely to the elimination of content confounds.

% ============================================================
\subsection{Things We Got Wrong}
\label{sec:things_wrong}

Transparency about error is not merely a rhetorical strategy; it is a methodological necessity in a field where the risk of self-deception about one's own results is non-trivial---particularly when studying deception.
We catalog the findings from our earlier work that were later shown to be artifacts or errors:

\begin{enumerate}[leftmargin=2em]
  \item \textbf{Bloom taxonomy inverted-U.}
  The inverted-U relationship between cognitive complexity and geometric richness was 90--98\% token count confound.
  Higher-complexity Bloom categories (evaluate, create) elicit longer responses, which produce mechanically richer geometry.
  The pattern vanishes under \fwl residualization.

  \item \textbf{RDCT naming.}
  The RDCT paradigm was labeled as ``cache truncation'' but was implemented as prompt perturbation.
  The results (geometric sensitivity to prompt variation) remain valid, but the claimed connection to Watson's $1/e$ falsification criterion \cite{watson2024}---which requires genuine truncation---was invalid.

  \item \textbf{Identity deduplication.}
  Within-prompt identity similarity of 100\% was reported as evidence for identity encoding stability.
  The 100\% figure was an artifact of non-deduplicated runs: the same trial appeared multiple times in the dataset.
  After deduplication, cross-prompt similarity (92--97\%) remains a genuine finding, but the within-prompt claim is retracted.

  \item \textbf{Deception expansion/compression split.}
  An early summary claimed that some deception models showed geometric compression.
  All seven models showed expansion; the claimed split was a transcription error that propagated through several documents before the audit caught it.

  \item \textbf{Per-layer uniformity (Hackathon Exp~35).}
  During the hackathon, initial analysis reported uniform deception signal ($d > 1.0$) across all 28 layers.
  Token-controlled reanalysis via ANCOVA revealed that this uniformity was a token-count artifact.
  The corrected finding: per-token Cohen's $d$ \emph{reverses} to $-1.909$, meaning deceptive generation is actually \emph{sparser per token} at every layer.
  The ANCOVA controls confirm this at 27/28 layers ($p < 0.05$).

  \item \textbf{Sycophancy detection (Hackathon Exp~39).}
  The original sycophancy experiment used 16 trials per condition and reported \auroc $= 0.9375$ for honest-vs-sycophantic detection.
  An independent input-length audit \cite{dwayne_audit} demonstrated that a classifier trained on input token counts alone achieved \auroc $= 0.948$, exceeding the geometric classifier.
  The sycophancy signal was entirely an input-length artifact.
  We replaced this experiment with a rigorous same-prompt design (Section~\ref{sec:sycophancy}): 50 prompts $\times$ 4 conditions $\times$ 3 models, system prompt spread $\leq 3$ tokens, yielding \fwl \auroc of 0.824--1.000 for honest-vs-sycophantic detection and 0.757--0.942 for the stricter sycophantic-vs-contrarian comparison.
  The replacement experiment confirms that the sycophancy signal is real but was inflated in the original by approximately 0.11 \auroc points from the input-length confound.
\end{enumerate}

We report these errors not as a weakness but as evidence that the surviving results have been stress-tested.
Every major finding reported in Sections~\ref{sec:results}--\ref{sec:invariances} has survived: (a)~the Kavi audit's 39 independent statistical tests, (b)~Dwayne's independent falsification pipeline \cite{dwayne_audit} (22 test files, 300+ tests, which identified the Exp~39 input-length confound and 2 other falsified comparisons), (c)~the Phase~3d seven-test red-team battery (6/7 pass, with the one failure well-characterized), (d)~the mandatory \fwl residualization that eliminates the most common confound class, and (e)~the same-prompt design that eliminates content confounds.
The methodology that discovered these errors---systematic adversarial testing, external audit, and mandatory confound control---is itself a contribution of this work, and we encourage its adoption in the broader interpretability community.
